%%=====================================================================
%%  Демо-отчёт: показывает все возможности шаблона.
%%  Новый отчёт: make new DIR=../lab-2  (заготовка — src/template.tex)
%%
%%  ПРАВИЛО ОДНО: блочная команда (\suaitasks, \suailist, \suaitable,
%%  \suaieq, \suaisources) читает строки под собой до ПУСТОЙ строки.
%%=====================================================================
\documentclass[a4paper,14pt]{extarticle}
\usepackage{suai-report}

%%  Титульный лист. Новые отчёты берут эти данные отсюда
%%  или из соседнего отчёта.
\suaisetup{
  department   = 41,
  teacher-post = Ассистент,
  teacher      = Д. С. Лисюков,
  type         = Отчет о лабораторной работе,  % прописными — автоматически
  number       = 1,                            % пусто — без номера
  title        = Знакомство с менеджером заметок Obsidian,
  course       = Облачные технологии и сервисы,
  group        = 4414,
  student      = Г. В. Моисеев,                % несколько: А. А. Иванов, Б. Б. Петров
  date         = 12.09.2026,                   % today — сегодня; пусто — от руки
% course-label = по дисциплине:,
% year         = 2026,                         % по умолчанию текущий
}

\begin{document}

\maketitle
\suaitoc

%%---------------------------------------------------------------------
\suaiintro

Цель работы — изучить менеджер заметок Obsidian и настроить
синхронизацию хранилища между устройствами через облачный сервис.

\suaitasks
установить Obsidian и создать хранилище
освоить разметку Markdown и связи между заметками
настроить синхронизацию через облачное хранилище
оценить удобство работы с базой знаний

%%---------------------------------------------------------------------
\section{Ход работы}

\subsection{Установка и настройка}

Obsidian установлен с официального сайта. Окно программы после
создания хранилища показано на~\figref{scheme}.

\suaiimg{scheme}{Главное окно Obsidian с открытым хранилищем}

Основные элементы интерфейса:
\suailist
панель файлов слева
редактор заметок в центре
граф связей и панель ссылок справа

\subsection{Сравнение способов синхронизации}

Способы синхронизации сравниваются в~\tabref{sync}.

\suaitable[sync]{Сравнение способов синхронизации хранилища}
Способ          | Стоимость      | Шифрование | Скорость
Obsidian Sync   | 4 \$ в месяц   | да         | высокая
Яндекс Диск     | бесплатно      | нет        | средняя
Git-репозиторий | бесплатно      | нет        | низкая

\subsection{Оценка объёма хранилища}

Объём хранилища оценивается по~\formref{size}:
\suaieq[size]{V = N \cdot \bar{s}}
V | объём хранилища, КБ
N | число заметок
\bar{s} | средний размер заметки, КБ

Для 300 заметок по 4 КБ получаем 1,2 МБ.

\subsection{Скрипт резервного копирования}

Скрипт для резервной копии хранилища приведён в~\lstref{backup}.

\begin{code}[bash, backup]{Резервное копирование хранилища}
#!/bin/bash
tar -czf "vault-$(date +%F).tar.gz" ~/Obsidian/Vault
\end{code}

%%---------------------------------------------------------------------
\suaiconclusion

В ходе работы установлен Obsidian, создано хранилище и настроена
синхронизация. Для личного использования достаточно бесплатного
облачного диска, для команды удобнее Obsidian Sync.

%%---------------------------------------------------------------------
\suaisources
Obsidian Help [Электронный ресурс]. URL: https://help.obsidian.md (дата обращения: 12.09.2026).
Фаулер М. Архитектура корпоративных программных приложений. М.: Вильямс, 2015. 544 с.

\end{document}
